\documentclass[12pt]{article}

\usepackage[T2A]{fontenc}
\usepackage[utf8]{inputenc}
\usepackage[russian]{babel}
\usepackage{amsmath}
\usepackage{amssymb}
\usepackage{graphicx}
\usepackage{cite}

\title{Исследование применения искусственного интеллекта для генерации LaTeX}
\author{Салимов Б. Г.}
\date{\today}

\begin{document}

\maketitle

\begin{abstract}
В данной работе мы исследуем агентское программное обеспечение для помощи в сборке научной статьи.
\end{abstract}

\noindent\textbf{Ключевые слова:} ИИ, LaTeX

\section{Введение}
Давно многие ученые хотели заниматься только исследованиями и созданием новых идей, но чтобы статьи за них оформлял кто-то. Кто знает, может этот момент настал?

В качестве примера приведём ссылку на статью о статистическом анализе ионосферной модели~\cite{Basciftci2023}.

В таблице~\ref{tab:stars} приведены характеристики первых пятнадцати звёзд выборки: температура, светимость, радиус и абсолютная звёздная величина. Данные охватывают звёзды с температурой от 1939 до 3628~K, среди которых преобладают холодные звёзды с низкой светимостью и слабой абсолютной звёздной величиной.

\input{../cache/tables/stars}

\section{Современные автомобили}
Современный автомобильный мир переживает настоящий расцвет: производители по всему миру выпускают машины, которые сочетают передовые технологии, невероятную мощность и выразительный дизайн. Электрический Tesla Model S Plaid разгоняется до сотни менее чем за две секунды, а гиперкары вроде Bugatti Chiron и Rimac Nevera преодолевают отметку в 1500 лошадиных сил, тогда как Porsche 911 GT3 RS и Ferrari SF90 Stradale доказывают, что и классические бензиновые, и гибридные решения остаются на пике инженерии. Не отстают и массовые бренды: Toyota, BMW, Mercedes-Benz и Hyundai предлагают доступные модели с богатым оснащением, автопилотами и гибридными силовыми установками. Особое место занимают доступные спортивные автомобили — Mazda MX-5 и Toyota GR86, которые за умеренную цену дарят чистый драйверский азарт.

На рисунке~\ref{fig:engine} представлена зависимость мощности автомобиля от объёма его двигателя: в целом больший объём двигателя соответствует большей мощности, однако разброс значений показывает, что современные технологии позволяют получать высокую отдачу и с компактными моторами.

\begin{figure}[htbp]
\centering
\includegraphics[width=0.8\textwidth]{../cache/plot/engine_size_vs_horsepower.png}
\caption{Зависимость мощности от объёма двигателя}
\label{fig:engine}
\end{figure}

\bibliographystyle{unsrt}
\bibliography{references}

\end{document}
